% from ega1-3-fr.tex
\oldpage[I]{110}
D'où la proposition en raison de l'unicité de la solution d'un problème
d'application universelle.

Ce résultat s'exprime encore en écrivant l'égalité (à un isomorphisme
canonique près) $(X_{(S')})_{(S'')}=X_{(S'')}$, si aucune confusion n'est à
craindre, ou encore
\begin{equation}
  (X\times_S S')\times_{S'}S''=X\times_S S''~;
  \tag{3.3.9.1}\label{I.3.3.9.1-fr}
\end{equation}
le caractère fonctoriel de l'isomorphisme défini dans
(\hyperref[I.3.3.9-fr]{3.3.9}) s'exprime de même par la formule de
\emph{transitivité des images réciproques de morphismes}
\begin{equation}
  (f_{(S')})_{(S'')}=f_{(S'')}
  \tag{3.3.9.2}\label{I.3.3.9.2-fr}
\end{equation}
pour tout $S$-morphisme $f:X\to Y$.

\begin{corollary}[3.3.10]
\label{I.3.3.10-fr}
Si $X$ et $Y$ sont deux $S$-préschémas, il existe un isomorphisme canonique
fonctoriel du $S'$-préschéma $X_{(S')}\times_{S'}Y_{(S')}$ sur le
$S'$-préschéma $(X\times_S Y)_{(S')}$.
\end{corollary}

En effet, on a, à des isomorphismes canoniques près
\[
  (X\times_S S')\times_{S'}(Y\times_S S')
  =X\times_S(Y\times_S S')=(X\times_S Y)\times_S S'
\]
en vertu de (\hyperref[I.3.3.9.1-fr]{3.3.9.1}) et de l'associativité des
produits de $S$-préschémas.

Le caractère fonctoriel de l'isomorphisme défini dans
(\hyperref[I.3.3.10-fr]{3.3.10}) s'exprime par la formule
\begin{equation}
  (u_{(S')},v_{(S')})_{S'}=((u,v)_S)_{(S')}
  \tag{3.3.10.1}\label{I.3.3.10.1-fr}
\end{equation}
pour tout couple de $S$-morphismes $u:T\to X$, $v:T\to Y$.

En d'autres termes, le foncteur image réciproque $X_{(S')}$
\emph{commute à la formation des produits}~; il commute aussi à la formation
des sommes (\hyperref[I.3.2.8-fr]{3.2.8}).

\begin{corollary}[3.3.11]
\label{I.3.3.11-fr}
Soient $Y$ un $S$-préschéma, $f:X\to Y$ un morphisme qui fait de $X$ un
$Y$-préschéma (et par suite aussi un $S$-préschéma). Le préschéma $X_{(S')}$
s'identifie alors au produit $X\times_Y Y_{(S')}$, la projection
$X\times_Y Y_{(S')}\to Y_{(S')}$ s'identifiant à $f_{(S')}$.
\end{corollary}

Soit $\psi:Y\to S$ le morphisme structural de $Y$~; on a le diagramme
commutatif
\phantomsection\label{I.3.3.11.diagram-fr}
\[
  \xymatrix{
    S'\ar[d] &
    Y_{(S')}\ar[l]_{\psi_{(S')}}\ar[d] &
    X_{(S')}\ar[l]_{f_{(S')}}\ar[d]\\
    S & Y\ar[l]^{\psi} & X\ar[l]^{f}
  }
\]
Or $Y_{(S')}$ s'identifie à $S'_{(\psi)}$, et $X_{(S')}$ à
$S'_{(\psi\circ f)}$~; tenant compte de (\hyperref[I.3.3.9-fr]{3.3.9}) et de
(\hyperref[I.3.3.4-fr]{3.3.4}), on en déduit le corollaire.

\begin{env}[3.3.12]
\label{I.3.3.12-fr}
Soient $f:X\to X'$, $g:Y\to Y'$ deux $S$-morphismes qui sont des
\emph{monomorphismes} de préschémas (T, I, 1.1)~; alors $f\times_S g$ est un
\emph{monomorphisme}. En effet, si $p$ et $q$ sont les projections de
$X\times_S Y$, $p'$, $q'$ celles de $X'\times_S Y'$, $u$, $v$ deux
$S$-morphismes $T\to X\times_S Y$, la relation
$(f\times_S g)\circ u=(f\times_S g)\circ v$ entraîne
$p'\circ(f\times_S g)\circ u=p'\circ(f\times_S g)\circ v$, autrement dit
$f\circ p\circ u=f\circ p\circ v$, et comme $f$ est un monomorphisme,
$p\circ u=p\circ v$~; utilisant le fait que $g$ est un monomorphisme, on
obtient de même $q\circ u=q\circ v$, d'où $u=v$.
